\chapter{Definición del problema}
\label{cap:problema}

El comercio exterior chileno opera mayoritariamente bajo la modalidad de \textbf{carga
consolidada}: las empresas importadoras pequeñas y medianas ---una zapatería, una juguetería, una
ferretería--- no generan volumen suficiente para llenar un contenedor completo, por lo que su
mercancía viaja agrupada con la de otros clientes. Toda la carga se declara en un
\textbf{manifiesto electrónico} en formato XML, que la nave transmite al Servicio Nacional de
Aduanas antes de recalar. Como el puerto de Valparaíso no dispone de espacio para abrir
contenedores dentro de su recinto, esta labor se realiza en \textbf{terminales extraportuarios}
(puertos secos), donde ocurre el \textbf{desconsolidado}: se abre el contenedor, se extrae la
mercancía y se separa por cliente destinatario, para que cada empresa retire lo suyo.

El desconsolidado es el punto del proceso donde la custodia de la mercancía cambia de manos varias
veces en pocas horas y donde, históricamente, la información se ha registrado de la forma más
precaria. Esto genera problemas concretos:

\begin{itemize}
    \item \textbf{Registro manual en papel y planillas sueltas:} los datos se anotan a mano y se
          transcriben después, con errores, retrasos y sin posibilidad de consultar el estado de
          una faena en tiempo real.
    \item \textbf{Evidencia fotográfica dispersa:} las fotografías se toman con teléfonos
          personales y circulan por mensajería instantánea, por lo que son difíciles de localizar
          y de auditar semanas después.
    \item \textbf{Ausencia de prueba del estado de la carga al abrir el contenedor:} ante un
          reclamo, el terminal no puede demostrar si el daño venía desde origen o se produjo
          durante sus propias maniobras.
    \item \textbf{Indefinición de responsabilidades ante los seguros:} sin un registro fechado del
          sello y de la mercancía, no es posible determinar rápidamente a qué aseguradora
          corresponde responder.
    \item \textbf{Descoordinación en la planificación de faenas y cuadrillas:} genera tiempos
          muertos, sobreasignación de personal e incumplimiento de horarios comprometidos.
    \item \textbf{Reporte al cliente lento y heterogéneo:} el importador se entera del estado de su
          carga días después de la faena.
\end{itemize}

A lo anterior se suman restricciones del entorno físico: los patios están ocupados por miles de
contenedores metálicos que bloquean la señal Wi-Fi, cada faena genera decenas o cientos de
fotografías de alta resolución, el registro lo realiza un operario con guantes y a la intemperie, y
el terminal necesita registrar todo pero no necesariamente publicar todo antes de revisarlo.

Por lo tanto, el problema central es \textbf{cómo digitalizar y documentar el proceso de
desconsolidado, generando evidencia verificable del estado de la carga en cada etapa crítica, bajo
conectividad intermitente y sin aumentar la carga de trabajo del operario en terreno}.
